\setcounter{section}{23}

  \zwischenueberschrift{Soal latihan}

\inputaufgabe
{}
{

Bahaslah
rumus Newton--Leibniz
sebagai kasus khusus dari
teorema Stokes.

}
{} {}

\inputaufgabe
{}
{

Buktikan secara langsung
versi balok dari teorema Stokes
untuk balok
\mathl{[a_1,b_1] \times \cdots \times [a_n,b_n]}{} dan suatu
$n-1$-\definitionsverweis {bentuk diferensial}{}{}
$\omega$ berbentuk

\mavergleichskettedisp
{\vergleichskette
{ \omega
}
{ =} { x_1^{m_1 } \cdots x_n^{m_n }  dx_2 \wedge \ldots \wedge dx_n
}
{ } {
}
{ } {
}
{ } {
}
}
{}{}{.}

}
{} {}

\inputaufgabe
{}
{

Misalkan $M$ suatu
$n$-\definitionsverweis {dimensional}{}{}
\definitionsverweis {berorientasi}{}{}
\definitionsverweis {manifold diferensiabel berbatas}{}{}
dengan
\definitionsverweis {basis topologi terhitung}{}{},
dan misalkan $\omega$ suatu $(n-1)$-\definitionsverweis {bentuk diferensial}{}{}
yang
\definitionsverweis {terdiferensialkan secara kontinu}{}{},
\definitionsverweis {tertutup}{}{}, dan mempunyai
\definitionsverweis {tumpuan kompak}{}{}
pada $M$. Tunjukkan bahwa

\mavergleichskettedisp
{\vergleichskette
{ \int_{\partial M} \omega
}
{ =} { 0
}
{ } {
}
{ } {
}
{ } {
}
}
{}{}{.}

}
{} {}

\inputaufgabe
{}
{

Misalkan $M$ suatu
\definitionsverweis {manifold diferensiabel}{}{}
$n$-\definitionsverweis {dimensional}{}{},
\definitionsverweis {kompak}{}{}, dan
\definitionsverweis {berorientasi}{}{}
\zusatzklammer {tanpa batas} {} {},
dengan
\definitionsverweis {basis topologi terhitung}{}{},
dan misalkan $\omega$ suatu
$(n-1)$-\definitionsverweis {bentuk diferensial}{}{}
yang
\definitionsverweis {terdiferensialkan secara kontinu}{}{}
pada $M$. Tunjukkan bahwa

\mavergleichskettedisp
{\vergleichskette
{
\int_{ M  }  d\omega
}
{ =} { 0
}
{ } {
}
{ } {
}
{ } {
}
}
{}{}{.}

}
{Apa arti pernyataan ini untuk $S^1$? Bagaimana pernyataan tersebut dapat dibuktikan dalam kasus ini melalui suatu integral lintasan?} {}

\inputaufgabe
{}
{

Misalkan $M$ suatu
\definitionsverweis {manifold diferensiabel}{}{}
\definitionsverweis {kompak}{}{} dan
\definitionsverweis {berorientasi}{}{}
\zusatzklammer {tanpa batas} {} {},
dengan
\definitionsverweis {basis topologi terhitung}{}{},
dan misalkan $\tau$ suatu
\definitionsverweis {bentuk volume positif}{}{}
pada $M$. Tunjukkan bahwa $\tau$ tidak
\definitionsverweis {eksak}{}{.}

}
{Bagaimana keadaannya tanpa asumsi kekompakan?} {}

\inputaufgabegibtloesung
{}
{

Misalkan $M$ suatu
\definitionsverweis {manifold diferensiabel}{}{}
dan $\omega$ suatu
\definitionsverweis {eksak}{}{} \definitionsverweis {bentuk diferensial}{}{}
pada $M$. Misalkan $S$ suatu
\definitionsverweis {manifold diferensiabel}{}{}
\definitionsverweis {kompak}{}{} dan
\definitionsverweis {berorientasi}{}{}
\zusatzklammer {tanpa batas} {} {},
dengan
\definitionsverweis {basis topologi terhitung}{}{},
serta misalkan
\maabbdisp {\varphi} { S } { M
} {}
suatu pemetaan yang
\definitionsverweis {terdiferensialkan secara kontinu}{}{.}
Tunjukkan bahwa

\mavergleichskettedisp
{\vergleichskette
{
\int_{  S  }   \varphi^*\omega
}
{ =} { 0
}
{ } {
}
{ } {
}
{ } {
}
}
{}{}{.}

}
{} {}

\inputaufgabe
{}
{

Kita tinjau
\definitionsverweis {pemetaan yang terdiferensialkan secara kontinu}{}{}
\zusatzklammer {\mathlk{n \geq 2}{}} {} {}
\maabbeledisp {\pi} {\R^n \setminus \{0\}} { S^{n-1}
} {(x_1 , \ldots , x_n) } { { \frac{   1  }{  \sqrt{x_1^2 + \cdots + x_n^2}  } } (x_1 , \ldots , x_n)
} {.}
Misalkan $\omega$ adalah
\definitionsverweis {bentuk volume kanonik}{}{}
pada $S^{n-1}$. Tunjukkan bahwa
\mathl{\pi^* \omega}{} pada $\R^n \setminus \{0\}$ merupakan
\definitionsverweis {tertutup}{}{},
tetapi tidak
\definitionsverweis {eksak}{}{}
$n-1$-\definitionsverweis {bentuk diferensial}{}{.}

}
{} {}

\inputaufgabe
{}
{

Misalkan

\mavergleichskette
{\vergleichskette
{ M
}
{ = }{ B  \left(  0 , 1 \right) \setminus \{(0,0)\}
}
{ \subset }{ \R^2
}
{    }{
}
{   }{
}
}
{}{}{}
dengan batas

\mavergleichskette
{\vergleichskette
{ \partial M
}
{ = }{ S^1
}
{  }{
}
{    }{
}
{   }{
}
}
{}{}{.}
Kita tinjau kedua
$1$-\definitionsverweis {bentuk diferensial}{}{}
yang
\definitionsverweis {terdiferensialkan secara kontinu}{}{}

\mathdisp {\omega= ydx-xdy \text{ dan } \tau= { \frac{   1  }{  x^2+y^2 } } \left(  ydx-xdy  \right)} {  }

pada $M$. Tunjukkan bahwa pembatasan kedua bentuk tersebut pada batas saling berimpit, dan khususnya

\mavergleichskette
{\vergleichskette
{ \int_{\partial M} \omega
}
{ =} { \int_{\partial M} \tau
}
{  }{
}
{    }{
}
{   }{
}
}
{}{}{}
berlaku. Bandingkan
\mathkor {} {\int_M d \omega} {dan} {\int_M d \tau} {.}

}
{} {}

\inputaufgabe
{}
{

Yakinkan diri bahwa
teorema Green
tidak menyatakan bahwa luas suatu daerah berbatas di $\R^2$ hanya bergantung pada panjang batasnya.

}
{} {}

\inputaufgabe
{}
{

Misalkan $D$ segitiga dengan titik sudut
$(0,2),\, (1,-1)$, dan $(-2,-1)$,
serta

\mavergleichskette
{\vergleichskette
{ \tau
}
{ = }{ x^2ydx \wedge dy
}
{  }{
}
{    }{
}
{   }{
}
}
{}{}{}
suatu
$2$-\definitionsverweis {bentuk diferensial}{}{} pada $D$. Temukan suatu
\definitionsverweis {bentuk primitif}{}{}
untuk $\tau$, lalu gunakan bentuk tersebut untuk menghitung
\mathl{\int_{ D  }  \tau}{} melalui suatu integral pada batas segitiga.

}
{} {}

\inputaufgabe
{}
{

Verifikasikan
teorema Green
melalui perhitungan eksplisit untuk persegi satuan

\mavergleichskette
{\vergleichskette
{ T
}
{ = }{ [0,1] \times [0,1]
}
{  }{
}
{    }{
}
{   }{
}
}
{}{}{}
dan
\definitionsverweis {bentuk-bentuk diferensial}{}{}

\mavergleichskettedisp
{\vergleichskette
{ \omega
}
{ =} { x^ay^b dx + x^cy^d dy
}
{ } {
}
{ } {
}
{ } {
}
}
{}{}{}
dengan

\mavergleichskette
{\vergleichskette
{ a,b,c,d
}
{ \in }{ \N
}
{  }{
}
{    }{
}
{   }{
}
}
{}{}{.}

}
{} {}

\inputaufgabe
{}
{

Verifikasikan
teorema Green
melalui perhitungan eksplisit untuk himpunan
\mathl{T=[-2,2] \times [-2,2] \setminus U \left(   0 , 1  \right)}{}
\zusatzklammer {yakni persegi berpusat di asal dengan panjang sisi $4$, tanpa lingkaran satuan terbuka} {} {}
dan
\definitionsverweis {bentuk diferensial}{}{}

\mavergleichskette
{\vergleichskette
{ \omega
}
{ = }{(x-y)dx + xy dy
}
{  }{
}
{    }{
}
{   }{
}
}
{}{}{.}

}
{} {}

\inputaufgabegibtloesung
{}
{

Buktikan teorema Green untuk segitiga $D$ dengan titik-titik sudut

\mavergleichskette
{\vergleichskette
{ P_1,P_2,P_3
}
{ \in }{ \R^2
}
{  }{
}
{    }{
}
{   }{
}
}
{}{}{}
dan untuk bentuk diferensial
\mathl{xdy}{.}

}
{} {}

\inputaufgabe
{}
{

Misalkan
\maabbdisp {f} { [a,b] } { \R_{+}
} {}
suatu
\definitionsverweis {fungsi yang terdiferensialkan secara kontinu}{}{.}
Kita pandang
\definitionsverweis {subgraf}{}{}
sebagai suatu
\definitionsverweis {manifold berbatas}{}{},
dengan batas yang terdiri atas grafik, interval dasar, dan kedua sisi tegak, tetapi tanpa keempat titik sudutnya. Hitung luas subgraf dengan menggunakan kedua
\definitionsverweis {bentuk diferensial}{}{}

\mathl{xdy}{} dan
\mathl{ydx}{}
pada batas.

}
{} {}

\inputaufgabe
{}
{

Tentukan volume
\definitionsverweis {bola satuan tertutup}{}{}
tiga dimensi melalui suatu integral permukaan yang sesuai pada
\definitionsverweis {permukaan bola satuan}{}{}
$S^2$.

}
{} {}

\inputaufgabegibtloesung
{}
{

Kita tinjau
$2$-\definitionsverweis {bentuk diferensial}{}{}

\mavergleichskettedisp
{\vergleichskette
{ \omega
}
{ =} { x dy \wedge dz -y dx \wedge dz + z dx \wedge dy
}
{ } {
}
{ } {
}
{ } {
}
}
{}{}{}
pada bola satuan
\mathl{B  \left(  0 , 1  \right)}{.}
\aufzaehlungdreiabc{Tunjukkan bahwa $d \omega$ adalah tiga kali bentuk volume standar pada bola tersebut.
}{Tunjukkan bahwa
\mathl{\omega\vert_{S^2}}{} adalah bentuk luas standar pada permukaan bola satuan.
}{Hitung luas permukaan bola dari
volume bola
dengan menggunakan
teorema Stokes.
}

}
{} {}

\inputaufgabegibtloesung
{}
{

Misalkan $\omega$ suatu
$n-1$-\definitionsverweis {bentuk diferensial}{}{}
yang terdiferensialkan secara kontinu pada suatu
\definitionsverweis {himpunan terbuka}{}{}

\mavergleichskette
{\vergleichskette
{ U
}
{ \subseteq }{ \R^n
}
{  }{
}
{    }{
}
{   }{
}
}
{}{}{},
dan misalkan
\definitionsverweis {turunan eksterior}{}{}
diberikan oleh

\mavergleichskettedisp
{\vergleichskette
{d \omega
}
{ =} { fdx_1 \wedge \ldots \wedge dx_n
}
{ } {
}
{ } {
}
{ } {
}
}
{}{}{}
dengan suatu fungsi
\maabb {f} { U } {\R
} {.}
Tunjukkan untuk

\mavergleichskette
{\vergleichskette
{ P
}
{ \in }{ U
}
{  }{
}
{    }{
}
{   }{
}
}
{}{}{}
kesamaan

\mavergleichskettedisp
{\vergleichskette
{ f(P)
}
{ =} { { \frac{   1  }{  \beta_n  } } \cdot \operatorname{lim}_{   \epsilon \rightarrow  0  } \, { \frac{   1  }{  \epsilon^n } } \int_{S^{n-1}(P, \epsilon)} \omega
}
{ } {
}
{ } {
}
{ } {
}
}
{}{}{,}
dengan $\beta_n$ menyatakan
volume
bola satuan $n$-dimensional.

}
{} {}

Untuk suatu kasus khusus dari pernyataan di atas, lihat
Soal 58.23 (Analisis (Osnabrück 2021--2023)).

\inputaufgabe
{}
{

Misalkan $M$ suatu
\definitionsverweis {manifold diferensiabel}{}{}
dengan batas tak kosong. Tunjukkan bahwa terdapat suatu
\definitionsverweis {pemetaan diferensiabel}{}{}
\maabb {} { M } { \partial M
} {}
ada.

}
{} {}

\inputaufgabe
{}
{

Misalkan

\mavergleichskette
{\vergleichskette
{ H
}
{ \subseteq }{ \R^n
}
{  }{
}
{    }{
}
{   }{
}
}
{}{}{}
\zusatzklammer {
\mavergleichskette
{\vergleichskette
{ n
}
{ \geq }{ 1
}
{  }{
}
{    }{
}
{   }{
}
}
{}{}{}} {} {}
suatu
\definitionsverweis {setengah ruang}{}{.}
Tunjukkan bahwa terdapat suatu
\definitionsverweis {pemetaan diferensiabel}{}{}
\maabbdisp {\varphi} { H } { \partial H
} {}
yang
\definitionsverweis {pembatasannya}{}{}
pada
\mathl{\partial H}{} adalah
\definitionsverweis {identitas}{}{.}

}
{Bagaimana keadaannya untuk

\mavergleichskette
{\vergleichskette
{ n
}
{ = }{ 0
}
{  }{
}
{    }{
}
{   }{
}
}
{}{}{}?} {}

\inputaufgabe
{}
{

Kita tinjau
\definitionsverweis {manifold berbatas}{}{}

\mavergleichskette
{\vergleichskette
{ M
}
{ = }{ \R^n \setminus  U \left(   0 , 1  \right)
}
{  }{
}
{    }{
}
{   }{
}
}
{}{}{.}
Tunjukkan bahwa terdapat suatu pemetaan yang terdiferensialkan secara kontinu dari $M$ ke batas
\mathl{S \left(   0 ,  1  \right)}{}
yang merupakan identitas pada batas tersebut.

}
{} {}

Untuk soal berikut, Anda dapat menggunakan
Soal 21.8.

\inputaufgabe
{}
{

Kita tinjau
\definitionsverweis {manifold berbatas}{}{}

\mavergleichskettedisp
{\vergleichskette
{ M
}
{ =} { B  \left(  0 , 1 \right) \setminus \{  \left(  1   , \,   0                   \right) \}
}
{ \subseteq} { \R^2
}
{ } {
}
{ } {
}
}
{}{}{.}
Tunjukkan bahwa terdapat suatu
\definitionsverweis {pemetaan yang terdiferensialkan secara kontinu}{}{}
\maabbdisp {\varphi} { M } { \partial M
} {}
yang merupakan identitas pada batas.

}
{} {}

\inputaufgabe
{}
{

Tunjukkan bahwa pada suatu
\definitionsverweis {anulus}{}{}
terdapat pemetaan
\definitionsverweis {bijektif}{}{} yang
\definitionsverweis {terdiferensialkan secara kontinu}{}{}
dan tidak mempunyai
\definitionsverweis {titik tetap}{}{.}

}
{} {}

\inputaufgabe
{}
{

Tunjukkan bahwa pemetaan

\mavergleichskettedisp
{\vergleichskette
{ \varphi(x)
}
{ =} { x + \left( - \left\langle  x  , h(x)  \right\rangle + \sqrt{1 + \left\langle  x  , h(x)  \right\rangle^2 - \Vert { x} \Vert^2  }\right) \cdot h(x)
}
{ } {
}
{ } {
}
{ } {
}
}
{}{}{}
\zusatzklammer {dengan

\mavergleichskette
{\vergleichskette
{ h(x)
}
{ = }{ { \frac{   \psi(x)-x }{  \Vert { \psi (x)-x} \Vert  } }
}
{  }{
}
{    }{
}
{   }{
}
}
{}{}{}} {} {}
yang digunakan dalam pembuktian
teorema titik tetap Brouwer
memetakan bola satuan ke permukaan bola satuan.

}
{} {}

  \zwischenueberschrift{Tugas untuk dikumpulkan}

\inputaufgabe
{4}
{

Misalkan $D$ segitiga dengan titik sudut
$(0,2),\, (1,-1)$, dan $(-2,-1)$,
serta
\mathl{\tau =( 3x^2y^5-x \sin  y ) dx \wedge dy}{} suatu
$2$-\definitionsverweis {bentuk diferensial}{}{} pada $D$. Temukan suatu
\definitionsverweis {bentuk primitif}{}{} untuk $\tau$, lalu gunakan bentuk tersebut untuk menghitung
\mathl{\int_{ D  }  \tau}{} melalui suatu integral pada batas segitiga.

}
{} {}

\inputaufgabe
{6}
{

Kita tinjau kubus

\mavergleichskettedisp
{\vergleichskette
{ Q
}
{ =} { [-1,1]^3
}
{ \subseteq} { \R^3
}
{ } {
}
{ } {
}
}
{}{}{}
dan
$2$-\definitionsverweis {bentuk diferensial}{}{}

\mavergleichskettedisp
{\vergleichskette
{ \omega
}
{ =} { dx \wedge dy +y dx \wedge dz +x^2y^2z^2 dy \wedge dz
}
{ } {
}
{ } {
}
{ } {
}
}
{}{}{.}
Hitung
\mathl{d \omega}{} dan kedua integral
\mathkor {} {\int_{  \partial Q  }   \omega} {dan} {\int_{  Q  }  d \omega} {}
\zusatzklammer {secara terpisah} {} {.}

}
{} {}

\inputaufgabe
{4}
{

Hitung luas
\definitionsverweis {cakram satuan tertutup}{}{}
melalui suatu
\definitionsverweis {integral lintasan}{}{}
yang sesuai.

}
{} {}

\inputaufgabe
{2}
{

Tunjukkan bahwa pada suatu
\definitionsverweis {torus}{}{}
terdapat pemetaan
\definitionsverweis {bijektif}{}{} yang
\definitionsverweis {terdiferensialkan secara kontinu}{}{}
dan tidak mempunyai
\definitionsverweis {titik tetap}{}{.}

}
{} {}

\inputaufgabe
{3}
{

Misalkan $B$ adalah
\definitionsverweis {cakram satuan tertutup}{}{}
dan $K$ busur setengah lingkaran atas. Tunjukkan bahwa terdapat suatu
\definitionsverweis {pemetaan diferensiabel}{}{}
\maabbdisp {\varphi} { B } { K
} {}
yang
\definitionsverweis {pembatasannya}{}{}
pada
\mathl{K}{} adalah
\definitionsverweis {identitas}{}{.}

}
{} {}

\inputaufgabe
{3}
{

Misalkan

\mavergleichskette
{\vergleichskette
{ v,w
}
{ \in }{ \R^n
}
{  }{
}
{    }{
}
{   }{
}
}
{}{}{}
dengan
\mathkor {} {\Vert { v } \Vert \leq 1} {dan} {\Vert { w } \Vert = 1} {.}
Tunjukkan bahwa terdapat

\mavergleichskette
{\vergleichskette
{ a
}
{ \in }{ \R
}
{  }{
}
{    }{
}
{   }{
}
}
{}{}{}
sedemikian sehingga

\mavergleichskette
{\vergleichskette
{ \Vert {v+aw} \Vert
}
{ = }{ 1
}
{  }{
}
{    }{
}
{   }{
}
}
{}{}{.}

}
{} {}

\inputaufgabe
{4}
{

Kita tinjau
\definitionsverweis {manifold berbatas}{}{}

\mavergleichskettedisp
{\vergleichskette
{ M
}
{ =} { B  \left(  0 , 1 \right) \setminus \{  \left(  1   , \,   0                   \right) \}
}
{ \subseteq} { \R^2
}
{ } {
}
{ } {
}
}
{}{}{.}
Tunjukkan bahwa terdapat suatu pemetaan bebas titik tetap
\definitionsverweis {yang terdiferensialkan secara kontinu}{}{}
\maabbdisp {\varphi} { M } { M
} {}
ada.

}
{} {}
